% The paper beside the series (B120, agreed 18-19 Nov 2026): research level, self-contained (all
% definitions and statements from the start; papers 1-3 cited in remarks only), its own Lean package
% (Lean/best_of_order/, namespace BestOfOrder, one bridge file to AcceleratedCollatz; exact statements only,
% the table of bests is script data), the failed detectors in Appendix A, one line each; no li/zeta analogy.
% For Zenodo, then linked from the OEIS sequence of the bests per Order, A402443. Unnumbered in the series.
% Draft (19 Nov 2026), Zenodo version 1: Orders 1..347 (B148); version 2 to 399 when boo_scan is DONE. \fillin marks values still to come (none now).
% Scripts: _keep/tree_01/ (boo.py, bnb.py, boo_scan.cpp, run_boo_scan.sh) and _keep/tree_01/zeta/.
\documentclass[10pt,a4paper,titlepage]{article}
\usepackage[utf8]{inputenc}
\usepackage[T1]{fontenc}
\usepackage{amsmath,amssymb,amsthm}
\usepackage[margin=1.25in]{geometry}
\usepackage{hyperref}
\hypersetup{colorlinks, allcolors=blue!50!black}
\usepackage{booktabs}
\usepackage{array}
\usepackage{longtable}
\usepackage{xcolor}
\usepackage{tikz}
\usepackage{etoolbox}
\IfFileExists{imakeidx.sty}
  {\usepackage{imakeidx}\makeindex[intoc]}
  {\usepackage{makeidx}\makeindex}

% Lean declaration
\newcommand{\lean}[1]{\textsf{#1}}
% Skeleton only: the plan of a proof, and a value still to fill in
\newcommand{\plan}[1]{\par\noindent{\small\color{black!60}\emph{Proof plan.} #1}\par}
\newcommand{\fillin}[1]{{\color{red!60!black}[#1]}}

% Layout: paper (default) or screen: pdflatex "\def\screenversion{}\input{file}" (as papers 1 and 2)
\newif\ifpaper
\papertrue
\ifdefined\screenversion \paperfalse \fi
% Paragraphs are never split across pages (paper version), so a chapter's first lines stay with its heading
\raggedbottom
\ifpaper
\interlinepenalty=10000
\postdisplaypenalty=10000
\fi
\ExplSyntaxOn
\seq_new:N \l__bw_parts_seq
\bool_new:N \l__bw_first_bool
\cs_new_protected:Npn \__bw_part:n #1
  {
    \fcolorbox{black!50}{black!5}
      {
        \bool_set_true:N \l__bw_first_bool
        \tl_map_inline:nn {#1}
          {
            \bool_if:NF \l__bw_first_bool { \textcolor{black!30}{\vrule width 0.3pt} }
            \bool_set_false:N \l__bw_first_bool
            \makebox[0.62em]{\rule[-0.4ex]{0pt}{2.2ex}\ttfamily ##1}
          }
      }
  }
% binary word: digits in cells, parts split at "-"
\NewDocumentCommand \bw { m }
  {
    \text
      {
        \group_begin:
        \setlength{\fboxsep}{0pt}
        \setlength{\fboxrule}{0.4pt}
        \seq_set_split:Nnn \l__bw_parts_seq { - } {#1}
        \seq_map_indexed_inline:Nn \l__bw_parts_seq
          {
            \int_compare:nNnF {##1} = {1} { \hspace{0.1em}\texttt{-}\hspace{0.1em} }
            \__bw_part:n {##2}
          }
        \group_end:
      }
  }
\ExplSyntaxOff
% Words without frames: pdflatex "\def\plainwords{}\input{file}"
\ifdefined\plainwords \RenewDocumentCommand \bw { m } {\text{\texttt{#1}}} \fi

% the admissible words of Orders 1 to 6 as one tree, the best of each Order bold, generated by
% Project_kb/tree_tikz.py (N = 6); the same tree, with options, is drop_tree.html
\newcommand{\droptreefig}{%
% generated by tree_tikz.py (N = 6)
\definecolor{treea}{HTML}{C0392B}
\definecolor{treeb}{HTML}{D35400}
\definecolor{treec}{HTML}{B7950B}
\definecolor{treed}{HTML}{1E8449}
\definecolor{treee}{HTML}{1F618D}
\definecolor{treef}{HTML}{7D3C98}
\node[font=\scriptsize, text=gray] at (0.000,0.450) {1};
\node[font=\scriptsize, text=gray] at (0.520,0.450) {2};
\node[font=\scriptsize, text=gray] at (1.040,0.450) {3};
\node[font=\scriptsize, text=gray] at (1.560,0.450) {4};
\node[font=\scriptsize, text=gray] at (2.080,0.450) {5};
\node[font=\scriptsize, text=gray] at (2.600,0.450) {6};
\node[font=\scriptsize, text=gray] at (3.120,0.450) {7};
\node[font=\scriptsize, text=gray] at (3.640,0.450) {8};
\node[font=\scriptsize, text=gray] at (4.160,0.450) {9};
\node[font=\scriptsize, text=gray] at (4.680,0.450) {10};
\draw[gray, line width=0.4pt] (0.000,0.000) -- (0.520,0.000);
\node[font=\tiny, text=gray, above] at (0.260,0.000) {0};
\draw[gray, line width=0.4pt] (0.000,0.000) -- (0.000,0.000) -- (0.520,-0.400);
\node[font=\tiny, text=gray, below left] at (0.260,-0.200) {1};
\draw[gray, line width=0.4pt] (0.520,-0.400) -- (1.040,-0.400);
\node[font=\tiny, text=gray, above] at (0.780,-0.400) {0};
\draw[gray, line width=0.4pt] (1.040,-0.400) -- (1.560,-0.400);
\node[font=\tiny, text=gray, above] at (1.300,-0.400) {0};
\draw[gray, line width=0.4pt] (1.040,-0.400) -- (1.040,-0.400) -- (1.560,-0.800);
\node[font=\tiny, text=gray, below left] at (1.300,-0.600) {1};
\draw[gray, line width=0.4pt] (1.560,-0.800) -- (2.080,-0.800);
\node[font=\tiny, text=gray, above] at (1.820,-0.800) {0};
\draw[gray, line width=0.4pt] (1.560,-0.800) -- (1.560,-0.800) -- (2.080,-1.200);
\node[font=\tiny, text=gray, below left] at (1.820,-1.000) {1};
\draw[gray, line width=0.4pt] (2.080,-1.200) -- (2.600,-1.200);
\node[font=\tiny, text=gray, above] at (2.340,-1.200) {0};
\draw[gray, line width=0.4pt] (2.600,-1.200) -- (3.120,-1.200);
\node[font=\tiny, text=gray, above] at (2.860,-1.200) {0};
\draw[gray, line width=0.4pt] (2.600,-1.200) -- (2.600,-1.200) -- (3.120,-1.600);
\node[font=\tiny, text=gray, below left] at (2.860,-1.400) {1};
\draw[gray, line width=0.4pt] (3.120,-1.600) -- (3.640,-1.600);
\node[font=\tiny, text=gray, above] at (3.380,-1.600) {0};
\draw[gray, line width=0.4pt] (3.120,-1.600) -- (3.120,-1.600) -- (3.640,-2.000);
\node[font=\tiny, text=gray, below left] at (3.380,-1.800) {1};
\draw[gray, line width=0.4pt] (3.640,-2.000) -- (4.160,-2.000);
\node[font=\tiny, text=gray, above] at (3.900,-2.000) {0};
\draw[gray, line width=0.4pt] (4.160,-2.000) -- (4.680,-2.000);
\node[font=\tiny, text=gray, above] at (4.420,-2.000) {0};
\draw[gray, line width=0.4pt] (2.080,-1.200) -- (2.080,-2.000) -- (2.600,-2.400);
\node[font=\tiny, text=gray, below left] at (2.340,-2.200) {1};
\draw[gray, line width=0.4pt] (2.600,-2.400) -- (3.120,-2.400);
\node[font=\tiny, text=gray, above] at (2.860,-2.400) {0};
\draw[gray, line width=0.4pt] (3.120,-2.400) -- (3.640,-2.400);
\node[font=\tiny, text=gray, above] at (3.380,-2.400) {0};
\draw[gray, line width=0.4pt] (3.120,-2.400) -- (3.120,-2.400) -- (3.640,-2.800);
\node[font=\tiny, text=gray, below left] at (3.380,-2.600) {1};
\draw[gray, line width=0.4pt] (3.640,-2.800) -- (4.160,-2.800);
\node[font=\tiny, text=gray, above] at (3.900,-2.800) {0};
\draw[gray, line width=0.4pt] (4.160,-2.800) -- (4.680,-2.800);
\node[font=\tiny, text=gray, above] at (4.420,-2.800) {0};
\draw[gray, line width=0.4pt] (2.600,-2.400) -- (2.600,-2.800) -- (3.120,-3.200);
\node[font=\tiny, text=gray, below left] at (2.860,-3.000) {1};
\draw[gray, line width=0.4pt] (3.120,-3.200) -- (3.640,-3.200);
\node[font=\tiny, text=gray, above] at (3.380,-3.200) {0};
\draw[gray, line width=0.4pt] (3.640,-3.200) -- (4.160,-3.200);
\node[font=\tiny, text=gray, above] at (3.900,-3.200) {0};
\draw[gray, line width=0.4pt] (4.160,-3.200) -- (4.680,-3.200);
\node[font=\tiny, text=gray, above] at (4.420,-3.200) {0};
\draw[gray, line width=0.4pt] (0.520,-0.400) -- (0.520,-3.200) -- (1.040,-3.600);
\node[font=\tiny, text=gray, below left] at (0.780,-3.400) {1};
\draw[gray, line width=0.4pt] (1.040,-3.600) -- (1.560,-3.600);
\node[font=\tiny, text=gray, above] at (1.300,-3.600) {0};
\draw[gray, line width=0.4pt] (1.560,-3.600) -- (2.080,-3.600);
\node[font=\tiny, text=gray, above] at (1.820,-3.600) {0};
\draw[gray, line width=0.4pt] (1.560,-3.600) -- (1.560,-3.600) -- (2.080,-4.000);
\node[font=\tiny, text=gray, below left] at (1.820,-3.800) {1};
\draw[gray, line width=0.4pt] (2.080,-4.000) -- (2.600,-4.000);
\node[font=\tiny, text=gray, above] at (2.340,-4.000) {0};
\draw[gray, line width=0.4pt] (2.600,-4.000) -- (3.120,-4.000);
\node[font=\tiny, text=gray, above] at (2.860,-4.000) {0};
\draw[gray, line width=0.4pt] (2.600,-4.000) -- (2.600,-4.000) -- (3.120,-4.400);
\node[font=\tiny, text=gray, below left] at (2.860,-4.200) {1};
\draw[gray, line width=0.4pt] (3.120,-4.400) -- (3.640,-4.400);
\node[font=\tiny, text=gray, above] at (3.380,-4.400) {0};
\draw[gray, line width=0.4pt] (3.120,-4.400) -- (3.120,-4.400) -- (3.640,-4.800);
\node[font=\tiny, text=gray, below left] at (3.380,-4.600) {1};
\draw[gray, line width=0.4pt] (3.640,-4.800) -- (4.160,-4.800);
\node[font=\tiny, text=gray, above] at (3.900,-4.800) {0};
\draw[gray, line width=0.4pt] (4.160,-4.800) -- (4.680,-4.800);
\node[font=\tiny, text=gray, above] at (4.420,-4.800) {0};
\draw[gray, line width=0.4pt] (2.080,-4.000) -- (2.080,-4.800) -- (2.600,-5.200);
\node[font=\tiny, text=gray, below left] at (2.340,-5.000) {1};
\draw[gray, line width=0.4pt] (2.600,-5.200) -- (3.120,-5.200);
\node[font=\tiny, text=gray, above] at (2.860,-5.200) {0};
\draw[gray, line width=0.4pt] (3.120,-5.200) -- (3.640,-5.200);
\node[font=\tiny, text=gray, above] at (3.380,-5.200) {0};
\draw[gray, line width=0.4pt] (3.120,-5.200) -- (3.120,-5.200) -- (3.640,-5.600);
\node[font=\tiny, text=gray, below left] at (3.380,-5.400) {1};
\draw[gray, line width=0.4pt] (3.640,-5.600) -- (4.160,-5.600);
\node[font=\tiny, text=gray, above] at (3.900,-5.600) {0};
\draw[gray, line width=0.4pt] (4.160,-5.600) -- (4.680,-5.600);
\node[font=\tiny, text=gray, above] at (4.420,-5.600) {0};
\draw[gray, line width=0.4pt] (2.600,-5.200) -- (2.600,-5.600) -- (3.120,-6.000);
\node[font=\tiny, text=gray, below left] at (2.860,-5.800) {1};
\draw[gray, line width=0.4pt] (3.120,-6.000) -- (3.640,-6.000);
\node[font=\tiny, text=gray, above] at (3.380,-6.000) {0};
\draw[gray, line width=0.4pt] (3.640,-6.000) -- (4.160,-6.000);
\node[font=\tiny, text=gray, above] at (3.900,-6.000) {0};
\draw[gray, line width=0.4pt] (4.160,-6.000) -- (4.680,-6.000);
\node[font=\tiny, text=gray, above] at (4.420,-6.000) {0};
\draw[gray, line width=0.4pt] (1.040,-3.600) -- (1.040,-6.000) -- (1.560,-6.400);
\node[font=\tiny, text=gray, below left] at (1.300,-6.200) {1};
\draw[gray, line width=0.4pt] (1.560,-6.400) -- (2.080,-6.400);
\node[font=\tiny, text=gray, above] at (1.820,-6.400) {0};
\draw[gray, line width=0.4pt] (2.080,-6.400) -- (2.600,-6.400);
\node[font=\tiny, text=gray, above] at (2.340,-6.400) {0};
\draw[gray, line width=0.4pt] (2.600,-6.400) -- (3.120,-6.400);
\node[font=\tiny, text=gray, above] at (2.860,-6.400) {0};
\draw[gray, line width=0.4pt] (2.600,-6.400) -- (2.600,-6.400) -- (3.120,-6.800);
\node[font=\tiny, text=gray, below left] at (2.860,-6.600) {1};
\draw[gray, line width=0.4pt] (3.120,-6.800) -- (3.640,-6.800);
\node[font=\tiny, text=gray, above] at (3.380,-6.800) {0};
\draw[gray, line width=0.4pt] (3.120,-6.800) -- (3.120,-6.800) -- (3.640,-7.200);
\node[font=\tiny, text=gray, below left] at (3.380,-7.000) {1};
\draw[gray, line width=0.4pt] (3.640,-7.200) -- (4.160,-7.200);
\node[font=\tiny, text=gray, above] at (3.900,-7.200) {0};
\draw[gray, line width=0.4pt] (4.160,-7.200) -- (4.680,-7.200);
\node[font=\tiny, text=gray, above] at (4.420,-7.200) {0};
\draw[gray, line width=0.4pt] (2.080,-6.400) -- (2.080,-7.200) -- (2.600,-7.600);
\node[font=\tiny, text=gray, below left] at (2.340,-7.400) {1};
\draw[gray, line width=0.4pt] (2.600,-7.600) -- (3.120,-7.600);
\node[font=\tiny, text=gray, above] at (2.860,-7.600) {0};
\draw[gray, line width=0.4pt] (3.120,-7.600) -- (3.640,-7.600);
\node[font=\tiny, text=gray, above] at (3.380,-7.600) {0};
\draw[gray, line width=0.4pt] (3.120,-7.600) -- (3.120,-7.600) -- (3.640,-8.000);
\node[font=\tiny, text=gray, below left] at (3.380,-7.800) {1};
\draw[gray, line width=0.4pt] (3.640,-8.000) -- (4.160,-8.000);
\node[font=\tiny, text=gray, above] at (3.900,-8.000) {0};
\draw[gray, line width=0.4pt] (4.160,-8.000) -- (4.680,-8.000);
\node[font=\tiny, text=gray, above] at (4.420,-8.000) {0};
\draw[gray, line width=0.4pt] (2.600,-7.600) -- (2.600,-8.000) -- (3.120,-8.400);
\node[font=\tiny, text=gray, below left] at (2.860,-8.200) {1};
\draw[gray, line width=0.4pt] (3.120,-8.400) -- (3.640,-8.400);
\node[font=\tiny, text=gray, above] at (3.380,-8.400) {0};
\draw[gray, line width=0.4pt] (3.640,-8.400) -- (4.160,-8.400);
\node[font=\tiny, text=gray, above] at (3.900,-8.400) {0};
\draw[gray, line width=0.4pt] (4.160,-8.400) -- (4.680,-8.400);
\node[font=\tiny, text=gray, above] at (4.420,-8.400) {0};
\draw[gray, line width=0.4pt] (1.560,-6.400) -- (1.560,-8.400) -- (2.080,-8.800);
\node[font=\tiny, text=gray, below left] at (1.820,-8.600) {1};
\draw[gray, line width=0.4pt] (2.080,-8.800) -- (2.600,-8.800);
\node[font=\tiny, text=gray, above] at (2.340,-8.800) {0};
\draw[gray, line width=0.4pt] (2.600,-8.800) -- (3.120,-8.800);
\node[font=\tiny, text=gray, above] at (2.860,-8.800) {0};
\draw[gray, line width=0.4pt] (3.120,-8.800) -- (3.640,-8.800);
\node[font=\tiny, text=gray, above] at (3.380,-8.800) {0};
\draw[gray, line width=0.4pt] (3.120,-8.800) -- (3.120,-8.800) -- (3.640,-9.200);
\node[font=\tiny, text=gray, below left] at (3.380,-9.000) {1};
\draw[gray, line width=0.4pt] (3.640,-9.200) -- (4.160,-9.200);
\node[font=\tiny, text=gray, above] at (3.900,-9.200) {0};
\draw[gray, line width=0.4pt] (4.160,-9.200) -- (4.680,-9.200);
\node[font=\tiny, text=gray, above] at (4.420,-9.200) {0};
\draw[gray, line width=0.4pt] (2.600,-8.800) -- (2.600,-9.200) -- (3.120,-9.600);
\node[font=\tiny, text=gray, below left] at (2.860,-9.400) {1};
\draw[gray, line width=0.4pt] (3.120,-9.600) -- (3.640,-9.600);
\node[font=\tiny, text=gray, above] at (3.380,-9.600) {0};
\draw[gray, line width=0.4pt] (3.640,-9.600) -- (4.160,-9.600);
\node[font=\tiny, text=gray, above] at (3.900,-9.600) {0};
\draw[gray, line width=0.4pt] (4.160,-9.600) -- (4.680,-9.600);
\node[font=\tiny, text=gray, above] at (4.420,-9.600) {0};
\draw[gray, line width=0.4pt] (2.080,-8.800) -- (2.080,-9.600) -- (2.600,-10.000);
\node[font=\tiny, text=gray, below left] at (2.340,-9.800) {1};
\draw[gray, line width=0.4pt] (2.600,-10.000) -- (3.120,-10.000);
\node[font=\tiny, text=gray, above] at (2.860,-10.000) {0};
\draw[gray, line width=0.4pt] (3.120,-10.000) -- (3.640,-10.000);
\node[font=\tiny, text=gray, above] at (3.380,-10.000) {0};
\draw[gray, line width=0.4pt] (3.640,-10.000) -- (4.160,-10.000);
\node[font=\tiny, text=gray, above] at (3.900,-10.000) {0};
\draw[gray, line width=0.4pt] (4.160,-10.000) -- (4.680,-10.000);
\node[font=\tiny, text=gray, above] at (4.420,-10.000) {0};
\draw[treea, opacity=0.75, line width=3.6pt, line join=round] (0.520,0.000) -- (0.000,0.000);
\draw[treeb, opacity=0.75, line width=3.1pt, line join=round] (1.560,-0.400) -- (1.040,-0.400) -- (0.520,-0.400) -- (0.000,0.000) -- (0.000,0.000);
\draw[treec, opacity=0.75, line width=2.6pt, line join=round] (2.080,-0.800) -- (1.560,-0.800) -- (1.040,-0.400) -- (1.040,-0.400) -- (0.520,-0.400) -- (0.000,0.000) -- (0.000,0.000);
\draw[treed, opacity=0.75, line width=2.2pt, line join=round] (3.120,-4.000) -- (2.600,-4.000) -- (2.080,-4.000) -- (1.560,-3.600) -- (1.560,-3.600) -- (1.040,-3.600) -- (0.520,-3.200) -- (0.520,-0.400) -- (0.000,0.000) -- (0.000,0.000);
\draw[treee, opacity=0.75, line width=1.7pt, line join=round] (3.640,-5.200) -- (3.120,-5.200) -- (2.600,-5.200) -- (2.080,-4.800) -- (2.080,-4.000) -- (1.560,-3.600) -- (1.560,-3.600) -- (1.040,-3.600) -- (0.520,-3.200) -- (0.520,-0.400) -- (0.000,0.000) -- (0.000,0.000);
\draw[treef, opacity=0.75, line width=1.2pt, line join=round] (4.680,-2.800) -- (4.160,-2.800) -- (3.640,-2.800) -- (3.120,-2.400) -- (3.120,-2.400) -- (2.600,-2.400) -- (2.080,-2.000) -- (2.080,-1.200) -- (1.560,-0.800) -- (1.560,-0.800) -- (1.040,-0.400) -- (1.040,-0.400) -- (0.520,-0.400) -- (0.000,0.000) -- (0.000,0.000);
\fill[darkgray] (0.000,0.000) circle (0.04);
\fill[darkgray] (0.520,-0.400) circle (0.04);
\fill[darkgray] (1.040,-0.400) circle (0.04);
\fill[darkgray] (1.560,-0.800) circle (0.04);
\fill[darkgray] (2.080,-1.200) circle (0.04);
\fill[darkgray] (2.600,-1.200) circle (0.04);
\fill[darkgray] (3.120,-1.600) circle (0.04);
\fill[darkgray] (3.640,-2.000) circle (0.04);
\fill[darkgray] (4.160,-2.000) circle (0.04);
\fill[darkgray] (2.600,-2.400) circle (0.04);
\fill[darkgray] (3.120,-2.400) circle (0.04);
\fill[darkgray] (3.640,-2.800) circle (0.04);
\fill[darkgray] (4.160,-2.800) circle (0.04);
\fill[darkgray] (3.120,-3.200) circle (0.04);
\fill[darkgray] (3.640,-3.200) circle (0.04);
\fill[darkgray] (4.160,-3.200) circle (0.04);
\fill[darkgray] (1.040,-3.600) circle (0.04);
\fill[darkgray] (1.560,-3.600) circle (0.04);
\fill[darkgray] (2.080,-4.000) circle (0.04);
\fill[darkgray] (2.600,-4.000) circle (0.04);
\fill[darkgray] (3.120,-4.400) circle (0.04);
\fill[darkgray] (3.640,-4.800) circle (0.04);
\fill[darkgray] (4.160,-4.800) circle (0.04);
\fill[darkgray] (2.600,-5.200) circle (0.04);
\fill[darkgray] (3.120,-5.200) circle (0.04);
\fill[darkgray] (3.640,-5.600) circle (0.04);
\fill[darkgray] (4.160,-5.600) circle (0.04);
\fill[darkgray] (3.120,-6.000) circle (0.04);
\fill[darkgray] (3.640,-6.000) circle (0.04);
\fill[darkgray] (4.160,-6.000) circle (0.04);
\fill[darkgray] (1.560,-6.400) circle (0.04);
\fill[darkgray] (2.080,-6.400) circle (0.04);
\fill[darkgray] (2.600,-6.400) circle (0.04);
\fill[darkgray] (3.120,-6.800) circle (0.04);
\fill[darkgray] (3.640,-7.200) circle (0.04);
\fill[darkgray] (4.160,-7.200) circle (0.04);
\fill[darkgray] (2.600,-7.600) circle (0.04);
\fill[darkgray] (3.120,-7.600) circle (0.04);
\fill[darkgray] (3.640,-8.000) circle (0.04);
\fill[darkgray] (4.160,-8.000) circle (0.04);
\fill[darkgray] (3.120,-8.400) circle (0.04);
\fill[darkgray] (3.640,-8.400) circle (0.04);
\fill[darkgray] (4.160,-8.400) circle (0.04);
\fill[darkgray] (2.080,-8.800) circle (0.04);
\fill[darkgray] (2.600,-8.800) circle (0.04);
\fill[darkgray] (3.120,-8.800) circle (0.04);
\fill[darkgray] (3.640,-9.200) circle (0.04);
\fill[darkgray] (4.160,-9.200) circle (0.04);
\fill[darkgray] (3.120,-9.600) circle (0.04);
\fill[darkgray] (3.640,-9.600) circle (0.04);
\fill[darkgray] (4.160,-9.600) circle (0.04);
\fill[darkgray] (2.600,-10.000) circle (0.04);
\fill[darkgray] (3.120,-10.000) circle (0.04);
\fill[darkgray] (3.640,-10.000) circle (0.04);
\fill[darkgray] (4.160,-10.000) circle (0.04);
\node[anchor=west, font=\scriptsize, text=gray] at (7.480,0.450) {class id};
\node[anchor=west, font=\scriptsize, text=gray] at (9.330,0.450) {best of Order};
\fill[treea] (0.520,0.000) circle (0.08);
\draw[treea, densely dotted, line width=0.3pt] (0.640,0.000) -- (5.080,0.000);
\node[anchor=west, font=\scriptsize, text=treea] at (5.130,0.000) {\textbf{1}\ \texttt{10}};
\node[anchor=east, font=\scriptsize, text=treea] at (8.630,0.000) {$1$};
\node[anchor=west, font=\scriptsize, text=treea] at (9.330,0.000) {$5 \to 4$, ratio $4/5$};
\fill[treeb] (1.560,-0.400) circle (0.08);
\draw[treeb, densely dotted, line width=0.3pt] (1.680,-0.400) -- (5.080,-0.400);
\node[anchor=west, font=\scriptsize, text=treeb] at (5.130,-0.400) {\textbf{2}\ \texttt{1100}};
\node[anchor=east, font=\scriptsize, text=treeb] at (8.630,-0.400) {$3$};
\node[anchor=west, font=\scriptsize, text=treeb] at (9.330,-0.400) {$3 \to 2$, ratio $2/3$};
\fill[treec] (2.080,-0.800) circle (0.08);
\draw[treec, densely dotted, line width=0.3pt] (2.200,-0.800) -- (5.080,-0.800);
\node[anchor=west, font=\scriptsize, text=treec] at (5.130,-0.800) {\textbf{3}\ \texttt{11010}};
\node[anchor=east, font=\scriptsize, text=treec] at (8.630,-0.800) {$11$};
\node[anchor=west, font=\scriptsize, text=treec] at (9.330,-0.800) {$11 \to 10$, ratio $10/11$};
\fill[treed] (3.120,-1.200) circle (0.08);
\draw[treed, densely dotted, line width=0.3pt] (3.240,-1.200) -- (5.080,-1.200);
\node[anchor=west, font=\scriptsize, text=treed] at (5.130,-1.200) {\textbf{4}\ \texttt{1101100}};
\node[anchor=east, font=\scriptsize, text=treed] at (8.630,-1.200) {$59$};
\fill[treee] (3.640,-1.600) circle (0.08);
\draw[treee, densely dotted, line width=0.3pt] (3.760,-1.600) -- (5.080,-1.600);
\node[anchor=west, font=\scriptsize, text=treee] at (5.130,-1.600) {\textbf{5}\ \texttt{11011010}};
\node[anchor=east, font=\scriptsize, text=treee] at (8.630,-1.600) {$123$};
\fill[treef] (4.680,-2.000) circle (0.08);
\draw[treef, densely dotted, line width=0.3pt] (4.800,-2.000) -- (5.080,-2.000);
\node[anchor=west, font=\scriptsize, text=treef] at (5.130,-2.000) {\textbf{6}\ \texttt{1101101100}};
\node[anchor=east, font=\scriptsize, text=treef] at (8.630,-2.000) {$507$};
\fill[treee] (3.640,-2.400) circle (0.08);
\draw[treee, densely dotted, line width=0.3pt] (3.760,-2.400) -- (5.080,-2.400);
\node[anchor=west, font=\scriptsize, text=treee] at (5.130,-2.400) {\textbf{5}\ \texttt{11011100}};
\node[anchor=east, font=\scriptsize, text=treee] at (8.630,-2.400) {$219$};
\fill[treef] (4.680,-2.800) circle (0.08);
\draw[treef, densely dotted, line width=0.3pt] (4.800,-2.800) -- (5.080,-2.800);
\node[anchor=west, font=\scriptsize, text=treef] at (5.130,-2.800) {\textbf{6}\ \texttt{1101110100}};
\node[anchor=east, font=\scriptsize, text=treef] at (8.630,-2.800) {$347$};
\node[anchor=west, font=\scriptsize, text=treef] at (9.330,-2.800) {$347 \to 248$, ratio $248/347$};
\fill[treef] (4.680,-3.200) circle (0.08);
\draw[treef, densely dotted, line width=0.3pt] (4.800,-3.200) -- (5.080,-3.200);
\node[anchor=west, font=\scriptsize, text=treef] at (5.130,-3.200) {\textbf{6}\ \texttt{1101111000}};
\node[anchor=east, font=\scriptsize, text=treef] at (8.630,-3.200) {$923$};
\fill[treec] (2.080,-3.600) circle (0.08);
\draw[treec, densely dotted, line width=0.3pt] (2.200,-3.600) -- (5.080,-3.600);
\node[anchor=west, font=\scriptsize, text=treec] at (5.130,-3.600) {\textbf{3}\ \texttt{11100}};
\node[anchor=east, font=\scriptsize, text=treec] at (8.630,-3.600) {$23$};
\fill[treed] (3.120,-4.000) circle (0.08);
\draw[treed, densely dotted, line width=0.3pt] (3.240,-4.000) -- (5.080,-4.000);
\node[anchor=west, font=\scriptsize, text=treed] at (5.130,-4.000) {\textbf{4}\ \texttt{1110100}};
\node[anchor=east, font=\scriptsize, text=treed] at (8.630,-4.000) {$7$};
\node[anchor=west, font=\scriptsize, text=treed] at (9.330,-4.000) {$7 \to 5$, ratio $5/7$};
\fill[treee] (3.640,-4.400) circle (0.08);
\draw[treee, densely dotted, line width=0.3pt] (3.760,-4.400) -- (5.080,-4.400);
\node[anchor=west, font=\scriptsize, text=treee] at (5.130,-4.400) {\textbf{5}\ \texttt{11101010}};
\node[anchor=east, font=\scriptsize, text=treee] at (8.630,-4.400) {$199$};
\fill[treef] (4.680,-4.800) circle (0.08);
\draw[treef, densely dotted, line width=0.3pt] (4.800,-4.800) -- (5.080,-4.800);
\node[anchor=west, font=\scriptsize, text=treef] at (5.130,-4.800) {\textbf{6}\ \texttt{1110101100}};
\node[anchor=east, font=\scriptsize, text=treef] at (8.630,-4.800) {$583$};
\fill[treee] (3.640,-5.200) circle (0.08);
\draw[treee, densely dotted, line width=0.3pt] (3.760,-5.200) -- (5.080,-5.200);
\node[anchor=west, font=\scriptsize, text=treee] at (5.130,-5.200) {\textbf{5}\ \texttt{11101100}};
\node[anchor=east, font=\scriptsize, text=treee] at (8.630,-5.200) {$39$};
\node[anchor=west, font=\scriptsize, text=treee] at (9.330,-5.200) {$39 \to 38$, ratio $38/39$};
\fill[treef] (4.680,-5.600) circle (0.08);
\draw[treef, densely dotted, line width=0.3pt] (4.800,-5.600) -- (5.080,-5.600);
\node[anchor=west, font=\scriptsize, text=treef] at (5.130,-5.600) {\textbf{6}\ \texttt{1110110100}};
\node[anchor=east, font=\scriptsize, text=treef] at (8.630,-5.600) {$423$};
\fill[treef] (4.680,-6.000) circle (0.08);
\draw[treef, densely dotted, line width=0.3pt] (4.800,-6.000) -- (5.080,-6.000);
\node[anchor=west, font=\scriptsize, text=treef] at (5.130,-6.000) {\textbf{6}\ \texttt{1110111000}};
\node[anchor=east, font=\scriptsize, text=treef] at (8.630,-6.000) {$999$};
\fill[treed] (3.120,-6.400) circle (0.08);
\draw[treed, densely dotted, line width=0.3pt] (3.240,-6.400) -- (5.080,-6.400);
\node[anchor=west, font=\scriptsize, text=treed] at (5.130,-6.400) {\textbf{4}\ \texttt{1111000}};
\node[anchor=east, font=\scriptsize, text=treed] at (8.630,-6.400) {$15$};
\fill[treee] (3.640,-6.800) circle (0.08);
\draw[treee, densely dotted, line width=0.3pt] (3.760,-6.800) -- (5.080,-6.800);
\node[anchor=west, font=\scriptsize, text=treee] at (5.130,-6.800) {\textbf{5}\ \texttt{11110010}};
\node[anchor=east, font=\scriptsize, text=treee] at (8.630,-6.800) {$79$};
\fill[treef] (4.680,-7.200) circle (0.08);
\draw[treef, densely dotted, line width=0.3pt] (4.800,-7.200) -- (5.080,-7.200);
\node[anchor=west, font=\scriptsize, text=treef] at (5.130,-7.200) {\textbf{6}\ \texttt{1111001100}};
\node[anchor=east, font=\scriptsize, text=treef] at (8.630,-7.200) {$975$};
\fill[treee] (3.640,-7.600) circle (0.08);
\draw[treee, densely dotted, line width=0.3pt] (3.760,-7.600) -- (5.080,-7.600);
\node[anchor=west, font=\scriptsize, text=treee] at (5.130,-7.600) {\textbf{5}\ \texttt{11110100}};
\node[anchor=east, font=\scriptsize, text=treee] at (8.630,-7.600) {$175$};
\fill[treef] (4.680,-8.000) circle (0.08);
\draw[treef, densely dotted, line width=0.3pt] (4.800,-8.000) -- (5.080,-8.000);
\node[anchor=west, font=\scriptsize, text=treef] at (5.130,-8.000) {\textbf{6}\ \texttt{1111010100}};
\node[anchor=east, font=\scriptsize, text=treef] at (8.630,-8.000) {$815$};
\fill[treef] (4.680,-8.400) circle (0.08);
\draw[treef, densely dotted, line width=0.3pt] (4.800,-8.400) -- (5.080,-8.400);
\node[anchor=west, font=\scriptsize, text=treef] at (5.130,-8.400) {\textbf{6}\ \texttt{1111011000}};
\node[anchor=east, font=\scriptsize, text=treef] at (8.630,-8.400) {$367$};
\fill[treee] (3.640,-8.800) circle (0.08);
\draw[treee, densely dotted, line width=0.3pt] (3.760,-8.800) -- (5.080,-8.800);
\node[anchor=west, font=\scriptsize, text=treee] at (5.130,-8.800) {\textbf{5}\ \texttt{11111000}};
\node[anchor=east, font=\scriptsize, text=treee] at (8.630,-8.800) {$95$};
\fill[treef] (4.680,-9.200) circle (0.08);
\draw[treef, densely dotted, line width=0.3pt] (4.800,-9.200) -- (5.080,-9.200);
\node[anchor=west, font=\scriptsize, text=treef] at (5.130,-9.200) {\textbf{6}\ \texttt{1111100100}};
\node[anchor=east, font=\scriptsize, text=treef] at (8.630,-9.200) {$735$};
\fill[treef] (4.680,-9.600) circle (0.08);
\draw[treef, densely dotted, line width=0.3pt] (4.800,-9.600) -- (5.080,-9.600);
\node[anchor=west, font=\scriptsize, text=treef] at (5.130,-9.600) {\textbf{6}\ \texttt{1111101000}};
\node[anchor=east, font=\scriptsize, text=treef] at (8.630,-9.600) {$287$};
\fill[treef] (4.680,-10.000) circle (0.08);
\draw[treef, densely dotted, line width=0.3pt] (4.800,-10.000) -- (5.080,-10.000);
\node[anchor=west, font=\scriptsize, text=treef] at (5.130,-10.000) {\textbf{6}\ \texttt{1111110000}};
\node[anchor=east, font=\scriptsize, text=treef] at (8.630,-10.000) {$575$};
}

\theoremstyle{plain}
\newtheorem{theorem}{Theorem}[section]
\newtheorem{lemma}[theorem]{Lemma}
\newtheorem{proposition}[theorem]{Proposition}
\newtheorem{corollary}[theorem]{Corollary}
\theoremstyle{definition}
\newtheorem{definition}[theorem]{Definition}
\newtheorem{observation}[theorem]{Observation}
\theoremstyle{remark}
\newtheorem{remark}[theorem]{Remark}
\newtheorem{mainthm}{Theorem}
\renewcommand{\themainthm}{\Alph{mainthm}}
\emergencystretch=3em

\title{\textbf{Collatz, but stop at the drop}\\[0.5ex]
  \Large The start words with the highest drop ratio}
\author{\textbf{Ruud H.G. van Tol}}
\date{\textbf{Draft}}

\begin{document}
\maketitle

\bigskip
\begin{center}\textbf{Abstract}\end{center}
\begin{quotation}\small\noindent
Under the map $T(x) = x/2$ for even $x$ and $(3x+1)/2$ for odd $x$, every odd start $x > 1$ that
falls below itself does so after $L = p_2(n)$ steps, $n$ of them odd, where $p_2(n)$ is the bit
length of $3^n$; $n$ is the Order of the drop. Among the class ids of Order $n$ we ask which one
falls the least: the start $c$ with the highest drop ratio $T^L(c)/c$, the \emph{best of the Order}.
We show that the best is the word whose cycle point $x^\ast = \rho/(2^L - 3^n)$ lies
closest to $c$, give a window outside which no start can win, split the search exactly into a
start side and a drop side that meet in one congruence modulo $3^w$, and compute the bests of
Orders $1$ to $347$, each with a certificate. The statements are formalized in Lean~4 with
Mathlib, without \texttt{sorry}; the table of bests is script data.

\medskip\noindent\textbf{MSC2020:} 11B83 (primary); 90C32, 68V15.\\
\textbf{Keywords:} Collatz problem, $3x+1$ problem, parity vectors, fractional programming,
Dinkelbach method, Lean~4.
\end{quotation}

\section*{Main results}
\addcontentsline{toc}{section}{Main results}

Words, Orders, class ids and the constant $\rho$ are defined in Chapter~\ref{sec:words}; for a word
of Order $n$ and length $L = p_2(n)$, $2^L\,T^L(x) = 3^n x + \rho$ for every $x$ of its class.

\begin{mainthm}[The cycle point; Chapter~\ref{sec:cycle}]
For a class id $c$ of Order $n$, with $x^\ast = \rho/(2^L - 3^n)$,
\[
1 - \frac{T^L(c)}{c} = \Bigl(1 - \frac{3^n}{2^L}\Bigr)\Bigl(1 - \frac{x^\ast}{c}\Bigr),
\]
so the best of the Order is the class id with the largest $\theta = x^\ast/c$; $c$ drops exactly when
$c > x^\ast$, and $x^\ast$ is the strict minimum of its own cycle.
\end{mainthm}

\begin{mainthm}[The window; Chapter~\ref{sec:search}]
If a class id of Order $n$ has $T^L(c)/c > b/a > 3^n/2^L$, then $c\,(2^L b - 3^n a) < R(n)\,a$, where
$R(n) = \sum_{m<n} 3^{n-1-m} 2^{p_2(m)-1}$ is the largest $\rho$ of the Order. Seeded with any
class id of Order $n$, a scan over the starts in increasing order is complete once it passes the
window of the best found; so is a best-first search over the word tree.
\end{mainthm}

\begin{mainthm}[The meeting split; Chapter~\ref{sec:dinkelbach}]
Cut a word of Order $n$ right after its $w$-th odd step, after $m$ letters, and let
$y = T^m(c)$. Then for every $\lambda$
\[
\rho - \lambda c = \rho_p\Bigl(3^{n-w} + \frac{\lambda}{3^w}\Bigr) + 2^m\Bigl(\sigma - \frac{\lambda\,y}{3^w}\Bigr),
\]
with $\rho_p$ the constant of the start part and $\sigma$ that of the drop part. The two parts are
tied only by $2^m y \equiv \rho_p \pmod{3^w}$, and the start parts are exactly the words of Order
$w$ without their final $0$'s.
\end{mainthm}

\begin{mainthm}[The bests; Chapter~\ref{sec:search}, Appendix~\ref{app:table}]
The bests of Orders $1$ to $347$ are computed and certified (OEIS A402443~\cite{OEIS_A402443}). For $285$ of the Orders $1$ to
$347$ the best is the smallest start of its Order (OEIS A122442); at the others the $\theta$ of the smallest start
is more than $1/4$ of the best's (smallest at Order $272$).
\end{mainthm}

\subsection*{What was known}
The formula for $T^L$, the residue classes of the parity words and the density of the starts that
drop are due to Terras~\cite{Terras1976} and Everett~\cite{Everett1977}; the largest $\rho$ of an
Order and its unique word are Terras's~\cite[Thm.~2.6]{Terras1976}. The parametric method for
fractional programs is Dinkelbach's~\cite{Dinkelbach1967}; accelerated versions are in
\cite{DadushKNV2021} and~\cite{ChenZhengXia2025}. The smallest start of each Order is OEIS
A122442~\cite{OEIS_A122442}. The admissible words, their tree and the names used here are those of
the series this paper stands beside~\cite{Paper1,Paper2,Paper3}. As far as we know, the
question of the best of the Order and Theorems~A to~D are new.

\ifpaper\clearpage\fi
\tableofcontents

%------------------------------------------------------------------------------------------------
\section{Introduction}
\label{sec:intro}
Run the map $T(x) = x/2$ for even $x$, $T(x) = (3x+1)/2$ for odd $x$, from a start $x > 1$, and
stop at the drop: the first step at which the value falls below $x$. Even starts drop at once; an odd
start rises first, and its parity word up to the drop is an \emph{admissible} word: $n$ odd steps in
$L = p_2(n)$ steps, $p_2(n)$ the bit length of $3^n$, with $3^j > 2^\ell$ after every earlier step.
This $n$ is the \emph{Order} of the drop. Each admissible word is the start of one residue class modulo
$2^L$, and its least member above $1$, the \emph{class id} $c$, stands for the whole class
(Chapter~\ref{sec:words}).

All starts of a class are mapped by one affine map, $2^L T^L(x) = 3^n x + \rho$, so the drop ratio
$T^L(x)/x = 3^n/2^L + \rho/(2^L x)$ is highest at the class id. Every start of Order $n$ falls at least
to $3^n/2^L$ of itself; the constant $\rho$ is what holds it up. This paper asks, per Order: which class
id falls the least? We call it the \emph{best of the Order}. Its ratio is the one closest to $1$, so it
is the hardest case of the statement that every class id drops at its word, Terras's coefficient
stopping time conjecture, in the series called the width identity.

The answer has a clean form (Theorem~\ref{thm:cycle}): each word has a rational cycle point
$x^\ast = \rho/(2^L - 3^n)$, the strict minimum of its own rational cycle, and the best is the class id
with the largest share $\theta = x^\ast/c$. Since $\rho$ is at most the constant $R(n)$ of the wall word
(Chapter~\ref{sec:wall}), a start can only win inside a window (Theorem~\ref{thm:window}); this turns
the question into two exact searches, a best-first search over the word tree and a scan over the starts
in increasing order, both seeded with the smallest start of the Order. As a fractional program the
question invites Dinkelbach's method; cutting a word after its $w$-th one splits $\rho - \lambda c$ into
a start side and a drop side, tied only by one congruence modulo $3^w$, and the start sides are the
words of Order $w$ (Chapter~\ref{sec:dinkelbach}). Chapter~\ref{sec:search} gives the searches and their
cost, Chapter~\ref{sec:observations} what the bests of Orders $1$ to $347$ show, and
Chapter~\ref{sec:lean} the formalization in Lean. Appendix~\ref{app:failed} lists the ways to guess the
best that did not work, and Appendix~\ref{app:table} the bests.

The Lean package, the scripts and the data (the bests with their certificates) are at
\url{https://github.com/druud/Collatz-but-stop-at-the-drop}, in \texttt{src/best\_of\_order/}.
The paper, the Lean proofs and the scripts were made with Claude (Anthropic) as a working
partner; Chapter~\ref{sec:lean} tells how.

\begin{figure}[htbp]
\centering
\begin{tikzpicture}\droptreefig\end{tikzpicture}
\caption{The admissible words of Orders $1$ to $6$ as one tree, the letter number on top. From
each node the $0$ goes straight on and the $1$ to the next free row. Each word ends at its Order
(bold) and its class id. Bold paths: the best of each Order, with $c \to T^L(c)$.}
\label{fig:tree}
\end{figure}

%------------------------------------------------------------------------------------------------
\ifpaper\clearpage\fi
\section{Words and Orders}
\label{sec:words}

\begin{definition}[The map, the word, the Order]
$T(x) = x/2$ for even $x$, $T(x) = (3x+1)/2$ for odd $x$. The \emph{word} of $x$ is its parity
sequence $e_0 e_1 e_2 \ldots$, $e_\ell = T^\ell(x) \bmod 2$. Let $p_2(n)$ be the bit length of $3^n$,
so $2^{p_2(n)-1} < 3^n < 2^{p_2(n)}$. A word of length $L$ with $n$ letters $1$ is
\emph{admissible} of \emph{Order} $n$ if $3^{j} > 2^{\ell}$ for each nonempty proper prefix, with $j$ ones
among its $\ell$ letters, and $3^n < 2^L$. Then $L = p_2(n)$ and the last letter is $0$: the prefix of
length $L-1$ has $3^{j} > 2^{L-1}$; were the last letter $1$, then $j = n - 1$ and
$3^n = 3 \cdot 3^{n-1} > 3 \cdot 2^{L-1} > 2^L$. So $j = n$ and $2^{L-1} < 3^n < 2^L$.
\end{definition}

\begin{lemma}[Terras]
\label{lem:affine}
Let $w$ be a word of length $L$ with $n$ ones, the $q$-th at letter $k_q$ ($0$-based). For every $x$
whose word starts with $w$,
\[
2^L\,T^L(x) = 3^n x + \rho, \qquad \rho = \sum_{q=1}^{n} 3^{\,n-q}\,2^{k_q},
\]
and these $x$ form one residue class modulo $2^L$.
\end{lemma}
\begin{proof}
Induction on $L$; for $L = 0$, $\rho = 0$ and the class is all of $\mathbb{Z}$. Let the statement hold for
$w$ and append a letter $e$. If $e = 0$, then $T^L(x)$ is even and
$2^{L+1}T^{L+1}(x) = 2^L T^L(x) = 3^n x + \rho$. If $e = 1$, then
$2^{L+1}T^{L+1}(x) = 3 \cdot 2^L T^L(x) + 2^L = 3^{n+1}x + 3\rho + 2^L$, and $3\rho + 2^L$ is the constant
with a new one at letter $k_{n+1} = L$. For the classes: writing $x = 2^L i + r$ in the class of $w$,
$T^L(x) = 3^n i + T^L(r)$, whose parity is that of $i + T^L(r)$; so each letter $e$ fixes $i \bmod 2$,
and $wx$ with $e$ appended is one class modulo $2^{L+1}$. \cite[§1]{Terras1976}.
\end{proof}

\begin{definition}[Class id, drop value]
The \emph{class id} of an admissible word $w$ of Order $n$ is the least $c > 1$ whose word starts with $w$
(for the residue $r < 2^L$ of the class: $c = r$ if $r > 1$, else $c = r + 2^L$; Order~$1$ gives $c = 5$).
Its \emph{drop value} is $d = T^L(c)$.
\end{definition}

\begin{lemma}
\label{lem:congruence}
$c \equiv -\rho\,3^{-n} = -\sum_{q=1}^n 2^{k_q}\,3^{-q} \pmod{2^L}$.
\end{lemma}
\begin{proof}
$3^n c + \rho = 2^L T^L(c) \equiv 0 \pmod{2^L}$ by Lemma~\ref{lem:affine}, and $3$ is invertible
modulo $2^L$.
\end{proof}

\begin{lemma}
\label{lem:member}
Every $x > 1$ of the class of $w$ has $T^L(x)/x \le T^L(c)/c$, with equality only for $x = c$.
\end{lemma}
\begin{proof}
By Lemma~\ref{lem:affine}, $T^L(x)/x = 3^n/2^L + \rho/(2^L x)$ with the same $\rho > 0$ for the whole
class, and $x \ge c$ by the choice of $c$.
\end{proof}

\begin{definition}[Best of the Order]
The \emph{best of Order $n$} is the class id $c$ of Order $n$ with the largest drop ratio
$T^L(c)/c$. By Lemma~\ref{lem:member} it is also the start $x > 1$ of Order $n$ with the largest ratio.
\end{definition}

\begin{remark}
\label{rem:residue}
The words of all Orders form one binary tree (Figure~\ref{fig:tree}); \cite[Ch.~4]{Paper3}.
For Order $1$ the residue is $1$ and $c = 5$; for $n \ge 2$ the residue of an admissible word is
more than $1$ (the word of $1$ is $\bw{1010}\ldots$, and $\bw{10}$ is not a prefix of a word of Order
$n \ge 2$), so $c = r < 2^L$.
\end{remark}

\begin{remark}[Ties]
\label{rem:ties}
Two class ids of one Order can share their drop ratio: the first case is Order $9$, where $21223$ and
$25691$ have the same $\rho/c$; Order $13$ has $16$ class ids with one ratio, and Order $16$ has
$56$ sets of equal ratios. A class id $c$ ties $c^\ast$ only if $\rho c^\ast = \rho^\ast c$, so $c$ is a multiple of
$u = c^\ast/\gcd(c^\ast, \rho^\ast)$; and a tie of the best lies in the window of
Theorem~\ref{thm:window} (with $\le$ for $<$), so $c = ku$ with $k \le W/u$. For Orders $1$ to $347$
that leaves $6906$ candidates (at most $1063$, at Order $262$), and none is a tie
(\texttt{boo\_tie\_mult.py}, exact): the best has no tie for Orders $1$ to $347$. Independently, a
depth-first search of the words with the window of the best finds none for Orders $1$ to $229$
(\texttt{boo\_ties.cpp}). That it never has one is open. Where we say \emph{the} best, a tie would mean the smallest of them.
\end{remark}

%------------------------------------------------------------------------------------------------
\ifpaper\clearpage\fi
\section{The wall}
\label{sec:wall}

\begin{definition}[The wall word]
The \emph{wall word} of Order $n$ puts its $q$-th one at letter $p_2(q-1) - 1$, the last place
that keeps the prefix before it admissible.
\end{definition}

\begin{theorem}[Terras]
\label{thm:wall}
For every admissible word of Order $n$, $k_q \le p_2(q-1) - 1$ for each $q$, hence
\[
\rho \le R(n) = \sum_{q=1}^{n} 3^{\,n-q}\,2^{p_2(q-1)-1},
\]
with equality only for the wall word.
\end{theorem}
\begin{proof}
$k_1 = 0 = p_2(0) - 1$. For $q \ge 2$ the prefix before the $q$-th one is nonempty and proper, with
$q-1$ ones in $k_q$ letters, so $2^{k_q} < 3^{q-1} < 2^{p_2(q-1)}$ and $k_q \le p_2(q-1) - 1$. Hence
$R(n) - \rho = \sum_q 3^{n-q}\bigl(2^{p_2(q-1)-1} - 2^{k_q}\bigr) \ge 0$, with equality only if every
$k_q$ is on the wall. The wall word is admissible: a nonempty proper prefix of length $\ell$ with $j$
ones has $\ell \le p_2(j) - 1$ (the next one, at $p_2(j)-1$, is not in it; for $j = n$, $\ell < p_2(n)$),
so $2^\ell \le 2^{p_2(j)-1} < 3^j$. \cite[Thm.~2.6]{Terras1976}; \cite[Thm.~A]{Paper3}.
\end{proof}

\begin{corollary}[The tail bound]
\label{cor:tail}
If a word of Order $n$ starts with a prefix that has $j$ ones and constant $\rho_j$, then
$\rho \le 3^{n-j}\rho_j + \sum_{q>j} 3^{\,n-q}\,2^{p_2(q-1)-1} = 3^{n-j}\rho_j + R(n) - 3^{n-j}R(j)$.
\end{corollary}
\begin{proof}
$\rho = 3^{n-j}\rho_j + \sum_{q>j} 3^{n-q}2^{k_q}$, and $k_q \le p_2(q-1)-1$ by Theorem~\ref{thm:wall}.
The terms $q \le j$ of $R(n)$ sum to $3^{n-j}R(j)$.
\end{proof}

\begin{lemma}[The start parts]
\label{lem:startparts}
Let $1 \le w < n$. The prefixes of the words of Order $n$ that end in their $w$-th one are exactly the
words of Order $w$ without their final $0$'s.
\end{lemma}
\begin{proof}
Let $p$ be such a prefix, of length $m$. As $w < n$ it is a nonempty proper prefix, so $3^w > 2^m$,
and so are its own nonempty prefixes. Append $0$'s up to the least length $m + k$ with
$3^w < 2^{m+k}$: every prefix on the way keeps $3^w > 2^{m+i}$, so this is a word of Order $w$, and $p$
is that word without its final $0$'s. Conversely let $u$ be a word of Order $w$ and $p$ the word $u$
without its final $0$'s, of length $m$, ending in the $w$-th one. Then $p$ is a proper prefix of $u$,
so $3^w > 2^m$. Append $n - w$ ones: each new prefix has $3^{w+i} > 2^{m+i}$. Then append $0$'s up to
the least length with $3^n < 2^\ell$. Every nonempty proper prefix of the result is admissible, so it
is a word of Order $n$, and $p$ is its prefix ending in its $w$-th one. Cf.~\cite[Ch.~4]{Paper3}.
\end{proof}

%------------------------------------------------------------------------------------------------
\section{The cycle point}
\label{sec:cycle}

\begin{definition}[Cycle point, share]
For a word of Order $n$ let $x^\ast = \rho/(2^L - 3^n)$, the rational fixed point of
$x \mapsto (3^n x + \rho)/2^L$, and for its class id $c$ let $\theta = x^\ast/c$.
\end{definition}

\begin{theorem}
\label{thm:cycle}
$1 - T^L(c)/c = (1 - 3^n/2^L)(1 - \theta)$. Hence the best of the Order is the class id with
the largest $\theta$, and $T^L(c) < c$ exactly when $c > x^\ast$.
\end{theorem}
\begin{proof}
By Lemma~\ref{lem:affine}, $1 - T^L(c)/c = 1 - 3^n/2^L - \rho/(2^L c)$, and
$(1 - 3^n/2^L)\,\theta = \frac{2^L - 3^n}{2^L}\cdot\frac{\rho}{(2^L - 3^n)\,c} = \rho/(2^L c)$. As
$1 - 3^n/2^L > 0$, the ratio $T^L(c)/c$ grows with $\theta$, and $T^L(c) < c$ exactly when $\theta < 1$.
\end{proof}

\begin{remark}
That every class id $c > 1$ of every Order satisfies $c > x^\ast$ is the width identity of the series
(Terras's coefficient stopping time conjecture~\cite[Conj.~2.9]{Terras1976}; \cite[Thm.~A]{Paper3}):
open in general, and the best of the Order is its hardest case. It holds for every Order up to
$10\,781\,273$~\cite{Paper2}, so for all Orders of this paper. Before letter $L$ every start of the word
rises: $T^\ell(x) = (3^{w_\ell}x + \rho_\ell)/2^\ell > x$ for $0 < \ell < L$. So each start $x > 1$ of an
Order $n \le 10\,781\,273$ first falls below itself exactly at step $L = p_2(n)$.
\end{remark}

\begin{proposition}[The cycle]
\label{prop:cyclemin}
Let $x^\ast_\ell = (3^{w_\ell} x^\ast + \rho_\ell)/2^\ell$ be the points of the rational cycle of the
word, $w_\ell$ and $\rho_\ell$ the ones and the constant of its first $\ell$ letters. Then
$x^\ast_\ell > x^\ast$ for $0 < \ell < L$: the cycle point is the strict minimum of its cycle.
\end{proposition}
\begin{proof}
$x^\ast_L = x^\ast$ by the choice of $x^\ast$, and $x^\ast_{\ell+1}$ is $x^\ast_\ell/2$ or
$(3x^\ast_\ell + 1)/2$ by the letter $e_\ell$ (as in Lemma~\ref{lem:affine}): the points form a cycle of
the two affine maps. For $0 < \ell < L$,
$x^\ast_\ell - x^\ast = \bigl((3^{w_\ell} - 2^\ell)x^\ast + \rho_\ell\bigr)/2^\ell > 0$, since
$3^{w_\ell} > 2^\ell$ for the nonempty proper prefixes, $x^\ast > 0$ and $\rho_\ell > 0$.
\end{proof}

\begin{lemma}[The split of the trajectory]
\label{lem:split}
For $0 \le \ell \le L$: $T^\ell(c) = x^\ast_\ell + (c - x^\ast)\,3^{w_\ell}/2^\ell$.
\end{lemma}
\begin{proof}
$T^\ell(c) = (3^{w_\ell}c + \rho_\ell)/2^\ell$ by Lemma~\ref{lem:affine}. Both sides are affine in $c$
with slope $3^{w_\ell}/2^\ell$ and take the value $x^\ast_\ell$ at $c = x^\ast$.
\end{proof}

\begin{proposition}[Prefix cycle points]
\label{prop:prefixneg}
The cycle point of the prefix of length $\ell$ as a word of its own,
$\rho_\ell/(2^\ell - 3^{w_\ell})$, is negative for every nonempty proper prefix.
\end{proposition}
\begin{proof}
$\rho_\ell > 0$ and $2^\ell < 3^{w_\ell}$.
\end{proof}

\begin{observation}[A ceiling]
Let $t_n$ be the least integer $x \ge 1$ with $(3x+1)^n < 2^L x^n$: the constant start of a word whose
odd steps all multiply by $3 + 1/x$. Then $x^\ast(\text{best}) \le x^\ast(\text{wall}) = R(n)/(2^L - 3^n)$,
and $x^\ast(\text{wall}) < t_n$ for $n \le 2000$ (exact, \texttt{ceiling\_check.py}). $t_n$ depends on $L/n$ only.
\end{observation}

%------------------------------------------------------------------------------------------------
\section{Dinkelbach and the meeting split}
\label{sec:dinkelbach}

By Theorem~\ref{thm:cycle} the best maximizes $\rho/c$ over the words of Order $n$: a fractional
program over a finite set.

\begin{lemma}[Dinkelbach~\cite{Dinkelbach1967}]
\label{lem:dinkelbach}
Let $\lambda^\ast$ be the largest $\rho/c$ and $F(\lambda) = \max(\rho - \lambda c)$ over the words of
Order $n$. Then $F$ is strictly decreasing, $F(\lambda^\ast) = 0$, and from any class id the step
$\lambda \mapsto \rho/c$ of an argmax of $\rho - \lambda c$ reaches $\lambda^\ast$ in finitely many steps.
\end{lemma}
\begin{proof}
$F$ is a maximum of finitely many strictly decreasing functions of $\lambda$ ($c > 0$), so it is strictly
decreasing; $\rho - \lambda^\ast c \le 0$ for every word, with equality at the best, so $F(\lambda^\ast) = 0$.
Let $\lambda = \rho_0/c_0 \le \lambda^\ast$ be the ratio of some class id. Then $F(\lambda) \ge 0$. If
$F(\lambda) = 0$, then $\lambda = \lambda^\ast$. Otherwise an argmax has $\rho - \lambda c > 0$, so its
$\rho/c > \lambda$: the steps strictly increase within a finite set, and stop at $\lambda^\ast$.
\end{proof}

\begin{lemma}[The walk]
\label{lem:walk}
Let $n \ge 2$. With $c = \sum_\ell t_\ell 2^\ell$ ($t_\ell$ the bits of $c < 2^L$) and $e_\ell$ the letters,
\[
\rho - \lambda c = \sum_{\ell=0}^{L-1} 2^\ell \delta_\ell,\qquad \delta_\ell = e_\ell\,3^{\,n-w_\ell} - \lambda\,t_\ell,
\]
$w_\ell$ the ones among the letters $0..\ell$.
\end{lemma}
\begin{proof}
A one at letter $\ell = k_q$ has $w_\ell = q$, so $\rho = \sum_\ell e_\ell\,3^{n-w_\ell}2^\ell$ by
Lemma~\ref{lem:affine}; $c < 2^L$ by Remark~\ref{rem:residue}.
\end{proof}

\begin{theorem}[The meeting split]
\label{thm:split}
Cut a word of Order $n$ after its $w$-th one, after $m$ letters; let $\rho_p$ be the constant of the start
part, $\sigma$ that of the drop part (its own letters counted from $0$) and $y = T^m(c)$. Then
$\rho = 3^{n-w}\rho_p + 2^m\sigma$, $c = (2^m y - \rho_p)/3^w$, and
\[
\rho - \lambda c = \rho_p\Bigl(3^{n-w} + \frac{\lambda}{3^w}\Bigr) + 2^m\Bigl(\sigma - \frac{\lambda\,y}{3^w}\Bigr).
\]
Conversely, let $p$ be a start part (a word of Order $w$ without its final $0$'s,
Lemma~\ref{lem:startparts}) and $v$ a drop part with $n - w$ ones. Then $pv$ is a word of Order $n$
exactly when each nonempty proper prefix of $v$, with $j$ ones in $\ell$ letters, has
$3^{w+j} > 2^{m+\ell}$ and $3^n < 2^{m+|v|}$. Its members $x > 0$ are then exactly the numbers
$(2^m y - \rho_p)/3^w$ with $y > \rho_p/2^m$ in the class of $v$ and $2^m y \equiv \rho_p \pmod{3^w}$;
the class id is the least of them above $1$.
\end{theorem}
\begin{proof}
The ones of $v$ sit at $k = m + k'$ and are the $(w+q')$-th of $pv$, so Lemma~\ref{lem:affine} gives
$\rho = 3^{n-w}\rho_p + \sum_{q'} 3^{(n-w)-q'}2^{m+k'} = 3^{n-w}\rho_p + 2^m\sigma$, and for the start
part $2^m y = 3^w c + \rho_p$. Substituting $c$ gives the identity. Admissibility of $pv$: the nonempty
prefixes of $pv$ inside $p$ are admissible as $p$ is a proper prefix of a word of Order $w$ with
$3^w > 2^m$, and the longer ones are the condition on $s$. For the members: if $x$ has the word $pv$,
then $y = T^m(x)$ has the word $v$ and $3^w x = 2^m y - \rho_p$. Conversely, for such a $y$ the number
$x = (2^m y - \rho_p)/3^w$ is a positive integer with $3^w x \equiv -\rho_p \pmod{2^m}$, so $x$ is in the
class of $p$ (Lemma~\ref{lem:congruence}, whose proof holds for any word), and
$T^m(x) = (3^w x + \rho_p)/2^m = y$ has the word $v$. The congruence is the only tie between the two
parts: $p$ alone fixes $\rho_p$, $m$, $w$; $v$ alone fixes $\sigma$ and the class of $y$.
Checked exactly by \texttt{mitm.py}.
\end{proof}

\begin{remark}[Both sides in turns]
By Lemma~\ref{lem:startparts} the start side at cut $w$ is the set of words of Order $w$, the same
for every $n > w$. The drop side grows backward from the last letter on residues $y \bmod 2^k$
($y' = 2y$ or $(2y-1)/3$). A word can be the best only if each part plus the other side's maximum
is $\ge 0$; with the true maxima, the drop side keeps $1$ to $3$ words (measured, \texttt{mitm4.py}).
\end{remark}

%------------------------------------------------------------------------------------------------
\section{The search}
\label{sec:search}

\begin{theorem}[The window]
\label{thm:window}
Let $b/a > 3^n/2^L$ be the ratio of some class id of Order $n$. A class id $c$ of Order $n$ with
$T^L(c)/c > b/a$ satisfies $c\,(2^L b - 3^n a) < R(n)\,a$. For a prefix with $j$ ones, constant $\rho_j$ and
residue $r$, every class id below it with a higher ratio satisfies
$r\,(2^L b - 3^n a) < \bigl(3^{n-j}\rho_j + \sum_{q>j} 3^{\,n-q} 2^{p_2(q-1)-1}\bigr)\,a$.
\end{theorem}
\begin{proof}
$T^L(c)/c > b/a$ means $a\,2^L T^L(c) > b\,2^L c$, that is $a(3^n c + \rho) > 2^L b\,c$, that is
$c\,(2^L b - 3^n a) < \rho\,a$. Now $\rho \le R(n)$ by Theorem~\ref{thm:wall}; below a prefix,
$\rho$ is at most the bound of Corollary~\ref{cor:tail}, and $c \ge r$ as $c \equiv r \pmod{2^\ell}$,
$c > 0$, $0 < r < 2^\ell$.
\end{proof}

\begin{corollary}[Two exact searches]
\begin{enumerate}
\item \emph{Best-first.} Pop the prefixes by residue, prune by Theorem~\ref{thm:window}, stop at the first
residue beyond the window of the best found (\texttt{boo.py}).
\item \emph{The scan.} Run over the starts $x$ in increasing order, keep per Order the highest
$T^L(x)/x$; Order $n$ is final once every $x$ below the window of its best is scanned
(\texttt{boo\_scan.cpp}).
\end{enumerate}
Both are exact for any seed; the seed only sets the first window.
\end{corollary}
\begin{proof}
Best-first: a prefix of residue $r$ has only descendants of residue at least $r$, and its class ids are
at least $r$, so a popped residue beyond the window of the best found bounds every node still in the
heap; a pruned node holds no class id with a higher ratio by Theorem~\ref{thm:window}; the window only
shrinks as the best improves. The scan computes the first step at which $x$ falls below itself, and
the ones before it; by the remark after Theorem~\ref{thm:cycle} that is the Order of $x$'s word for
all Orders here. So every class id $c$ of Order $n$ is met as a start of Order $n$, and every other start of its class is larger with a lower ratio (Lemma~\ref{lem:member}); so once
all $x$ below the window of the best found are scanned, Theorem~\ref{thm:window} leaves no class id that
could beat it.
\end{proof}

\begin{remark}[Seeds]
The smallest start of Order $n$, $s_0 = a(n+1)$ of A122442, is the best itself for $98$ of the Orders $1$
to $124$ (and for $285$ of the Orders $1$ to $347$). A guess from the previous Order (its $\theta$, halved until a word beats it) is exact as
well but costs about twice the nodes, and certified candidates built from the bests of earlier Orders
beat $s_0$ only at Orders $6$, $8$ and $107$ (Appendix~\ref{app:failed}).
\end{remark}

\begin{remark}[Cost]
For Orders $1$ to $124$ the best-first search seeded with $s_0$ visits $164\,801\,730$ nodes in all
(Order $124$: $15\,147\,961$); the scan finishes Orders $1$ to $124$ at $x < 2^{26}$. The scan for
Orders $1$ to $347$ covers every start below $80\,887\,088\,311\,813$, the window of Order $346$ and the
largest of these Orders (Theorem~\ref{thm:window}), on a wheel modulo $2^{28}$ that leaves $3\,524\,586$ residues; Orders
$348$ to $399$ need it up to $17\,216\,187\,145\,311\,549$. Without a seed, the scan finds the smallest start of each
Order on the way, and its first start of each Order is checked against A122442.
\end{remark}

%------------------------------------------------------------------------------------------------
\ifpaper\clearpage\fi
\section{Observations}
\label{sec:observations}
% Data only, measured; each with its script.

\begin{observation}[The smallest start]
$\theta(s_0)/\theta(\text{best}) > 1/4$ for all Orders $1$ to $347$, smallest at Order $272$ (between
$25/100$ and $26/100$), then Orders $256$, $91$, $268$ and $221$; equality of $s_0$ and the best at $285$ of them
(\texttt{boo\_obs.py}).
\end{observation}

\begin{observation}[Backward ranking]
At $\lambda = \lambda^\ast$ let $S_k = \sum_{\ell=L-k}^{L-1} 2^\ell \delta_\ell$ be the sum of the walk over the
last $k$ letters. The best is the strict maximum of $S_k$ over all words of its Order for every
$k \ge k_0$, with (Orders $2$ to $21$, all words, \texttt{dinkel\_suffix.py})
\[
\begin{array}{r|rrrrrrrrrrrrrrrrrrrr}
n & 2&3&4&5&6&7&8&9&10&11&12&13&14&15&16&17&18&19&20&21\\
L & 4&5&7&8&10&12&13&15&16&18&20&21&23&24&26&27&29&31&32&34\\
k_0 & 1&1&3&2&6&4&8&7&9&9&12&13&19&14&16&16&18&24&21&23
\end{array}
\]
For Orders $15$ to $21$, $L/2 < k_0 < 4L/5$; Order $14$ needs $19$ of its $23$ letters.
\end{observation}

\begin{observation}[Merging trajectories]
The drop values $d = T^L(\text{best})$ of Orders $1$ to $347$ have $219$ odd parts: $61$ is the odd
part for $10$ Orders, $12\,729\,682\,859$ for $9$, $619\,739$, $92\,485\,801$, $1\,064\,232\,949$ and
$30\,643\,528\,277$ for $5$ each, $7$ more for $4$ each, $18$ for $3$ each, $38$ for $2$ each, and $150$
for one Order each. Seven of the odd parts
($31$, $167$, $325$, $61$, $23$, $5$, $1$) lie on the trajectory of $27$, for $21$ Orders (\texttt{boo\_hubs.py}). The inverse trees of such hubs hold the bests only mildly
more often than other starts (Appendix~\ref{app:failed}).
\end{observation}

%------------------------------------------------------------------------------------------------
\section{The Lean formalization}
\label{sec:lean}
The package \texttt{best\_of\_order/} (namespace \lean{BestOfOrder}) states and proves the
results below in Lean~4 with Mathlib, without \texttt{sorry}. Its words (\lean{List Bool}, a one for an
odd step), the map $T$ (\lean{tStep}), the word of a start (\lean{tBits}), the constant $\rho$
(\lean{rho}), the Order (\lean{order}), admissibility (\lean{Admissible}) and $R(n)$ (\lean{rhoMax}) come
from the library of the series~\cite{Paper2} through one bridge file, \texttt{Bridge.lean}. Only exact
statements are formalized; the table of bests (Appendix~\ref{app:table}) is script data, and so are
the two searches and the observations.

\begin{center}\small
\begin{tabular}{@{}ll@{}}
\toprule
Statement & Lean \\
\midrule
Lemma~\ref{lem:affine} & \lean{affine}, \lean{mem\_class\_iff} \\
Lemma~\ref{lem:congruence} & \lean{congruence} (as $-\rho\,3^{-n}$) \\
Lemma~\ref{lem:member} & \lean{ratio\_lt\_of\_lt} \\
Theorem~\ref{thm:wall} & \lean{rho\_le\_wall}, \lean{rho\_eq\_wall\_iff}, \lean{wall\_eq\_sum} \\
Corollary~\ref{cor:tail} & \lean{tail\_bound} \\
Lemma~\ref{lem:startparts} & \lean{isKey\_of\_prefix}, \lean{exists\_of\_key} \\
Theorem~\ref{thm:cycle} (A) & \lean{one\_sub\_ratio}, \lean{ratio\_lt\_ratio\_iff}, \lean{drops\_iff} \\
Proposition~\ref{prop:cyclemin} & \lean{cyclePoint\_lt\_cyclePt}, \lean{cyclePt\_length} \\
Lemma~\ref{lem:split} & \lean{iterate\_eq\_split} \\
Proposition~\ref{prop:prefixneg} & \lean{prefix\_cyclePoint\_neg} \\
Lemma~\ref{lem:dinkelbach} & \lean{F\_strictAnti}, \lean{F\_lamStar}, \lean{eq\_lamStar\_of\_F\_eq\_zero}, \lean{step\_lt} \\
Lemma~\ref{lem:walk} & \lean{rho\_eq\_sum}, \lean{walk} \\
Theorem~\ref{thm:split} (C) & \lean{split\_identity}, \lean{mem\_split\_iff}, \lean{mem\_of\_meet}, \lean{exists\_start\_iff} \\
Theorem~\ref{thm:window} (B) & \lean{window}, \lean{window\_prefix} \\
Remark~\ref{rem:ties} & \lean{tie\_rho}, \lean{tie\_dvd} \\
\bottomrule
\end{tabular}
\end{center}

The width identity, on which the scan rests (the remark after Theorem~\ref{thm:cycle}), is
\lean{WidthIdentity} of that library; its range to Order $10\,781\,273$ uses a range check outside
Lean~\cite{Paper2}.

\subsection*{How this was made}

The Lean files, the scripts and this text were written in conversation with Claude (Anthropic):
ideas were tried numerically first, in exact integers, then stated and proved in Lean, compiled
locally, and every open end was logged for the next session; the long runs (the scan, the tie
search) ran on the author's computer. Some of it went better than planned: the ties, first
searched word by word, turned out to need only the multiples of one number below the window
(Remark~\ref{rem:ties}). The author checked the statements and takes responsibility for them.
The same workflow is open to any reader: the Lean proofs can be checked on any computer, and the
scripts recompute the table of bests and its certificates.

\appendix
%------------------------------------------------------------------------------------------------
\ifpaper\clearpage\fi
\section{Detectors that fail}
\label{app:failed}
One line each, with the reason (measured; scripts in \texttt{src/best\_of\_order/scripts/}).
\begin{itemize}\setlength{\itemsep}{0.2ex}
\item \emph{The scaled wall word as an ideal staircase:} every word starts and ends below it; the best's
  $2$ to $8$ crossings are not special.
\item \emph{Ideal lines and curves (diagonal, ruler, $\log$-shaped):} they measure area only; the best
  ranks anywhere ($3$ to $1337$).
\item \emph{Trailing $1$-bits of the start:} no relation with the best.
\item \emph{Inverse trees:} the inverse tree of $91$ holds $16$ of the $106$ bests in $(91, 2^{22})$,
  $15\%$, against $9\%$ of the odd starts in that range: a mild excess only (\texttt{inv\_hub.py},
  \texttt{inv\_grow.py}).
\item \emph{The backward branch and bound:} correct, but costs more nodes than forward (Order $22$:
  $96\,743$ primed against $1\,505$).
\item \emph{A guessed $\theta$ from the previous Order:} exact with halving, about twice the nodes of the
  seed $s_0$.
\item \emph{Candidates from earlier Orders} (one $1$ inserted per Order step, beams of $8$): beat $s_0$ at
  Orders $6$, $8$, $107$ only.
\end{itemize}

%------------------------------------------------------------------------------------------------
\section{The bests}
\label{app:table}
The best $c$ of Order $n$ and its drop value $T^L(c)$; $\ast$: $c$ differs from the smallest start
$a(n+1)$ of A122442. The column $c$ is OEIS A402443~\cite{OEIS_A402443} (with $a(0) = 2$ for Order $0$).
For Orders $125$ to $347$ only the rows marked $\ast$ are listed; all bests are in the b-file of
A402443 and in the repository (\texttt{src/best\_of\_order/data/}), with a certificate file (\texttt{best\_cert\_347.txt}: $n$, $c$, $T^L(c)$
and the window of Theorem~\ref{thm:window}).

{\footnotesize
\begin{longtable}{@{}rrrc@{\qquad}rrrc@{}}
\toprule
$n$ & $c$ & $T^L(c)$ & & $n$ & $c$ & $T^L(c)$ & \\
\midrule
\endhead
1 & 5 & 4 &  & 63 & 101375 & 91535 & $\ast$ \\
2 & 3 & 2 &  & 64 & 226767 & 153563 &  \\
3 & 11 & 10 &  & 65 & 105639 & 53654 & $\ast$ \\
4 & 7 & 5 &  & 66 & 10087 & 7688 &  \\
5 & 39 & 38 &  & 67 & 284063 & 162307 & $\ast$ \\
6 & 347 & 248 & $\ast$ & 68 & 67583 & 57925 &  \\
7 & 231 & 124 &  & 69 & 90111 & 57925 &  \\
8 & 207 & 167 & $\ast$ & 70 & 45055 & 43444 &  \\
9 & 127 & 77 &  & 71 & 126575 & 91535 &  \\
10 & 359 & 325 &  & 72 & 299259 & 162307 &  \\
11 & 511 & 346 &  & 73 & 112511 & 91535 & $\ast$ \\
12 & 239 & 122 &  & 74 & 366447 & 223589 & $\ast$ \\
13 & 159 & 122 &  & 75 & 64255 & 58810 &  \\
14 & 1115 & 637 & $\ast$ & 76 & 84383 & 57925 &  \\
15 & 283 & 244 &  & 77 & 57115 & 29405 &  \\
16 & 991 & 637 &  & 78 & 56255 & 43444 &  \\
17 & 251 & 244 &  & 79 & 37503 & 21722 &  \\
18 & 167 & 122 &  & 80 & 60975 & 52975 &  \\
19 & 111 & 61 &  & 81 & 45127 & 29405 &  \\
20 & 1695 & 1378 &  & 82 & 393967 & 385043 &  \\
21 & 1307 & 797 &  & 83 & 827503 & 606566 & $\ast$ \\
22 & 871 & 797 &  & 84 & 1759951 & 967537 &  \\
23 & 927 & 637 &  & 85 & 35655 & 29405 &  \\
24 & 671 & 346 &  & 86 & 434223 & 268556 &  \\
25 & 155 & 122 &  & 87 & 495687 & 459857 &  \\
26 & 103 & 61 &  & 88 & 665215 & 462844 &  \\
27 & 1639 & 1424 &  & 89 & 2452159 & 1279622 & $\ast$ \\
28 & 91 & 61 &  & 90 & 528895 & 413995 &  \\
29 & 3431 & 3349 &  & 91 & 1242471 & 729412 & $\ast$ \\
30 & 3399 & 2489 &  & 92 & 437247 & 385043 &  \\
31 & 2287 & 1256 &  & 93 & 5825855 & 3847667 & $\ast$ \\
32 & 71 & 61 &  & 94 & 432923 & 428885 &  \\
33 & 6395 & 3949 &  & 95 & 565247 & 419981 &  \\
34 & 47 & 46 &  & 96 & 288615 & 160832 &  \\
35 & 31 & 23 &  & 97 & 376831 & 314986 &  \\
36 & 2047 & 1067 &  & 98 & 2548479 & 1597663 &  \\
37 & 27 & 23 &  & 99 & 611455 & 574993 &  \\
38 & 1819 & 1067 &  & 100 & 608111 & 428885 &  \\
39 & 25407 & 22331 & $\ast$ & 101 & 1585403 & 838607 &  \\
40 & 6887 & 4541 &  & 102 & 543515 & 431245 & $\ast$ \\
41 & 4591 & 4541 &  & 103 & 270271 & 160832 &  \\
42 & 13439 & 9968 &  & 104 & 362343 & 323434 &  \\
43 & 6383 & 3551 &  & 105 & 401151 & 268556 &  \\
44 & 4255 & 3551 &  & 106 & 4937243 & 2478956 & $\ast$ \\
45 & 7963 & 4984 &  & 107 & 3291495 & 2478956 & $\ast$ \\
46 & 7527 & 7066 &  & 108 & 6721703 & 3796783 &  \\
47 & 12399 & 8729 &  & 109 & 381727 & 323434 &  \\
48 & 7279 & 3844 &  & 110 & 667375 & 424094 &  \\
49 & 1583 & 1256 &  & 111 & 626331 & 597017 &  \\
50 & 1055 & 628 &  & 112 & 1733791 & 1239478 & $\ast$ \\
51 & 703 & 628 &  & 113 & 1564063 & 838607 &  \\
52 & 15039 & 10049 &  & 114 & 1541147 & 1239478 &  \\
53 & 186239 & 93317 & $\ast$ & 115 & 1027431 & 619739 &  \\
54 & 41407 & 31123 &  & 116 & 1345383 & 1217288 & $\ast$ \\
55 & 115547 & 65134 & $\ast$ & 117 & 1991615 & 1351492 &  \\
56 & 84975 & 71852 & $\ast$ & 118 & 1327743 & 675746 &  \\
57 & 147151 & 93317 & $\ast$ & 119 & 7303711 & 5575759 &  \\
58 & 34239 & 32572 &  & 120 & 6255855 & 3581858 &  \\
59 & 200863 & 143300 & $\ast$ & 121 & 6688623 & 5744467 & $\ast$ \\
60 & 99007 & 52975 &  & 122 & 7849755 & 5056270 &  \\
61 & 127647 & 102449 & $\ast$ & 123 & 3137471 & 3031417 &  \\
62 & 236991 & 142655 & $\ast$ & 124 & 9294427 & 6735181 &  \\
\bottomrule
\end{longtable}
}

\medskip\noindent Orders $125$ to $347$, the rows marked $\ast$:
{\footnotesize
\begin{longtable}{@{}rrrc@{\qquad}rrrc@{}}
\toprule
$n$ & $c$ & $T^L(c)$ & & $n$ & $c$ & $T^L(c)$ & \\
\midrule
130 & 30997755 & 15991322 & $\ast$ & 220 & 3119162651 & 2519100964 & $\ast$ \\
132 & 39091199 & 22687444 & $\ast$ & 221 & 2136191463 & 1293924665 & $\ast$ \\
136 & 27401343 & 20127209 & $\ast$ & 223 & 1483334895 & 1010789105 & $\ast$ \\
148 & 66742363 & 49693357 & $\ast$ & 234 & 7032588283 & 6476798801 & $\ast$ \\
153 & 43254511 & 30569932 & $\ast$ & 244 & 18023934831 & 14956412627 & $\ast$ \\
161 & 354899823 & 200885648 & $\ast$ & 247 & 26690574503 & 18687428219 & $\ast$ \\
167 & 151639911 & 122212106 & $\ast$ & 252 & 36068070847 & 23970706696 & $\ast$ \\
168 & 93596391 & 56574545 & $\ast$ & 256 & 109571933119 & 92164207369 & $\ast$ \\
174 & 182938751 & 157443683 & $\ast$ & 259 & 68821108463 & 48842559905 & $\ast$ \\
175 & 620227615 & 400342673 & $\ast$ & 261 & 45880738975 & 36631919929 & $\ast$ \\
177 & 183238907 & 133061072 & $\ast$ & 268 & 4167011775 & 3552816982 & $\ast$ \\
183 & 306533191 & 158466836 & $\ast$ & 272 & 180760272511 & 97527287242 & $\ast$ \\
186 & 998412031 & 870992600 & $\ast$ & 284 & 380685108059 & 208196864200 & $\ast$ \\
187 & 1012039419 & 662160625 & $\ast$ & 320 & 871451633243 & 496371686726 & $\ast$ \\
200 & 604499583 & 601365250 & $\ast$ & 321 & 580967755495 & 496371686726 & $\ast$ \\
203 & 2836010779 & 2380476962 & $\ast$ & 332 & 1764823531367 & 1018943581204 & $\ast$ \\
208 & 1080905007 & 861211745 & $\ast$ & 333 & 1178426794975 & 1020569816743 & $\ast$ \\
212 & 6893644607 & 3475732561 & $\ast$ & 338 & 1901445212927 & 1563112498897 & $\ast$ \\
\bottomrule
\end{longtable}
}

\ifpaper\clearpage\fi
\begin{thebibliography}{99}
\bibitem{ChenZhengXia2025} H.~Chen, C.~Zheng and Y.~Xia, \emph{Accelerated Dinkelbach method},
  arXiv:2510.26257 (2025).
\bibitem{DadushKNV2021} D.~Dadush, Z.~K. Koh, B.~Natura and L.~A. V\'egh, \emph{An accelerated
  Newton--Dinkelbach method and its application to two variables per inequality systems},
  arXiv:2004.08634.
\bibitem{Dinkelbach1967} W.~Dinkelbach, \emph{On nonlinear fractional programming}, Management Sci.\
  \textbf{13} (1967), 492--498.
\bibitem{Everett1977} C.~J. Everett, \emph{Iteration of the number-theoretic function
  $f(2n) = n$, $f(2n+1) = 3n+2$}, Adv.\ Math.\ \textbf{25} (1977), 42--45.
\bibitem{OEIS_A122442} OEIS Foundation Inc., \emph{Sequence A122442},
  \url{https://oeis.org/A122442}.
\bibitem{OEIS_A402443} OEIS Foundation Inc., \emph{Sequence A402443},
  \url{https://oeis.org/A402443}.
\bibitem{Terras1976} R.~Terras, \emph{A stopping time problem on the positive
  integers}, Acta Arith.\ \textbf{30} (1976), 241--252.
\bibitem{Paper1} Ruud H.G. van Tol, \emph{Collatz, but stop at the drop; Now with less
  noise} (2026), \url{https://doi.org/10.5281/zenodo.23631044}.
\bibitem{Paper2} Ruud H.G. van Tol and Mike Winkler, \emph{Collatz, but stop at the drop;
  A game with two letters: the proofs} (2026), \url{https://doi.org/10.5281/zenodo.23631547}.
\bibitem{Paper3} Ruud H.G. van Tol, \emph{Collatz, but stop at the drop; The game on the grid of
  $2^i 3^j$}, draft (2026).
\end{thebibliography}

\printindex

\end{document}
